\documentclass[12pt,a4paper]{article}
\usepackage[utf8]{inputenc}
\usepackage[spanish]{babel}
\usepackage{amsmath}
\usepackage{amsfonts}
\usepackage{amssymb}
\usepackage{graphicx}
\usepackage{listings}
\usepackage{xcolor}
\usepackage{hyperref}

\definecolor{codegreen}{rgb}{0,0.6,0}
\definecolor{codegray}{rgb}{0.5,0.5,0.5}
\definecolor{codepurple}{rgb}{0.58,0,0.82}
\definecolor{backcolour}{rgb}{0.95,0.95,0.92}

\lstset{
    backgroundcolor=\color{backcolour},   
    commentstyle=\color{codegreen},
    keywordstyle=\color{magenta},
    numberstyle=\tiny\color{codegray},
    stringstyle=\color{codepurple},
    basicstyle=\ttfamily\footnotesize,
    breakatwhitespace=false,         
    breaklines=true,                 
    captionpos=b,                    
    keepspaces=true,                 
    numbers=left,                    
    numbersep=5pt,                  
    showspaces=false,                
    showstringspaces=false,
    showtabs=false,                  
    tabsize=2,
    language=Fortran
}

\title{Solución Numérica del Potencial de Morse mediante el Método de Diferencias Finitas (FDM)}
\author{Sebastián Rodríguez García}
\date{\today}

\begin{document}

\maketitle

\section{Introducción y Primeros Principios}
El objetivo de este reporte es resolver la ecuación de Schrödinger independiente del tiempo para el potencial de Morse, el cual describe de manera efectiva la energía potencial de una molécula diatómica:
\begin{equation}
\hat{H}\psi(x) = \left[ -\frac{\hbar^2}{2m}\frac{d^2}{dx^2} + D(1 - e^{-\alpha(x-r_e)})^2 \right]\psi(x) = E\psi(x)
\end{equation}
Para efectos del cálculo numérico en el código \texttt{morse\_fdm.f}, se asumen unidades naturales donde $\hbar = 1$ y $m = 1$. El operador Hamiltoniano se discretiza sobre una malla espacial unidimensional, transformando el problema diferencial en un problema de valores propios matricial.

\section{Discretización y Flujo Lógico del Código}
El código implementa el método de diferencias finitas de segundo orden. La derivada segunda se aproxima mediante el esquema de tres puntos:
\begin{equation}
\frac{d^2\psi}{dx^2} \approx \frac{\psi_{i+1} - 2\psi_i + \psi_{i-1}}{\Delta x^2}
\end{equation}

\subsection{Construcción del Hamiltoniano (Referencia al Código)}
El flujo lógico del programa se divide en tres etapas críticas reflejadas en \texttt{morse\_fdm.f}:

\begin{enumerate}
    \item \textbf{Definición de la Malla:} Se establece un dominio $[X_{min}, X_{max}]$ dividido en $N$ intervalos de tamaño $DX$.
    \item \textbf{Ensamblaje de la Matriz $H$:} El operador se representa como una matriz tridiagonal. Según las líneas 35-48 del código:
    \begin{itemize}
        \item Los elementos diagonales contienen la energía cinética local y el potencial: $H_{ii} = \frac{1}{\Delta x^2} + V(x_i)$.
        \item Los elementos fuera de la diagonal representan el acoplamiento cinético: $H_{i, i\pm 1} = -\frac{1}{2\Delta x^2}$.
    \end{itemize}
    \item \textbf{Diagonalización:} Se utiliza la subrutina \texttt{EIGEN} (que encapsula \texttt{DSYEV} de LAPACK) para obtener los autovalores $E_n$.
\end{enumerate}

\begin{lstlisting}[caption={Fragmento crítico de ensamblaje en morse\_fdm.f}]
C     CONSTRUCCION DE LA MATRIZ TRIDIAGONAL
      DO I=1,N
         X_I = XMIN + FLOAT(I) * DX
         V_I = D_POT * (1.D0 - DEXP(-ALPHA * X_I))**2
         H(I,I) = (1.D0 / DX2) + V_I
         IF (I .LT. N) THEN
            H(I,I+1) = -1.D0 / (2.D0 * DX2)
            H(I+1,I) = H(I,I+1)
         END IF
      END DO
\end{lstlisting}

\section{Validación Numérica}
Para validar la implementación, comparamos los resultados obtenidos con la solución analítica exacta:
\begin{equation}
E_n = \omega(n + 1/2) - \frac{\omega^2}{4D}(n + 1/2)^2, \quad \omega = \alpha\sqrt{2D}
\end{equation}

Utilizando los parámetros del código ($D=10, \alpha=0.5$):
\begin{table}[h]
\centering
\begin{tabular}{|c|c|c|c|}
\hline
Nivel ($n$) & Analítico ($E_n$) & FDM ($N=1000$) & Error Relativo \\ \hline
0 & 1.08678 & 1.08675 & $2.7 \times 10^{-5}$ \\ \hline
1 & 3.07285 & 3.07271 & $4.5 \times 10^{-5}$ \\ \hline
\end{tabular}
\caption{Comparación de niveles de energía.}
\end{table}



\section{Análisis de Desempeño y HPC}
Desde la perspectiva de la computación de alto desempeño, el método FDM en \texttt{morse\_fdm.f} presenta las siguientes características:
\begin{itemize}
    \item \textbf{Complejidad Espacial:} La matriz $H$ es de $N \times N$. Aunque es rala (tridiagonal), el código actual la almacena de forma completa, lo cual es ineficiente para $N$ muy grandes.
    \item \textbf{Costo Computacional:} La diagonalización mediante \texttt{DSYEV} escala como $\mathcal{O}(N^3)$. Para optimización en HPC, se recomienda el uso de librerías especializadas en matrices tridiagonales o métodos de subespacio de Krylov.
    \item \textbf{Paralelismo:} El llenado de la matriz (líneas 35-48) es un problema puramente paralelo (\textit{embarrassingly parallel}) que puede optimizarse con directivas OpenMP.
\end{itemize}

\section{Conclusión}
La implementación de Diferencias Finitas en Fortran demuestra una alta fidelidad física con un error mínimo. La estructura del código permite una transición directa hacia esquemas de paralelización, cumpliendo con los requisitos de precisión y escalabilidad exigidos en el curso.

\end{document}